\chapter{Time, and why everything is multiplied by delta time}

\section{The problem}

Here is a piece of code that looks completely reasonable and is wrong:

\begin{lstlisting}
if (IsKeyDown(KEY_D)) player.x += 5;      // WRONG
\end{lstlisting}

It says ``move five pixels per frame''. On a 60 Hz monitor that is 300 pixels a
second. On a 144 Hz monitor it is 720. On a machine struggling at 30 FPS it is
150. Your game is a different game on every computer, and if the frame rate
stutters, the physics stutter with it.

\section{The fix}

\lstinline|GetFrameTime()| returns how long the previous frame took, in seconds.
People universally call it \emph{delta time}, or \lstinline|dt|. At 60 FPS it is
about \lstinline|0.0166|; at 120 FPS about \lstinline|0.0083|.

\begin{lstlisting}
float dt = GetFrameTime();
float speed = 300.0f;                         // PIXELS PER SECOND

if (IsKeyDown(KEY_D)) player.x += speed*dt;   // RIGHT
\end{lstlisting}

At 60 FPS that adds $300 \times 0.0166 = 5$ pixels. At 120 FPS it adds
$300 \times 0.0083 = 2.5$ pixels, twice as often. Same distance per second, on
every machine.

\keyidea{Every speed you write down is per \emph{second}, and every time you
apply it you multiply by \lstinline|dt|. Not sometimes --- always. Movement,
rotation, timers, fading, reloading, animation. If a number changes over time,
it is multiplied by \lstinline|dt|.}

\section{Reading a movement line}

Once the habit sticks, your constants start documenting themselves:

\begin{lstlisting}
float walkSpeed = 4.0f;       // metres per second  -> a person walking
float runSpeed  = 8.0f;       // metres per second  -> a person running
float turnSpeed = 90.0f;      // degrees per second -> a quarter turn each second
float reload    = 1.5f;       // seconds
\end{lstlisting}

Compare that with ``0.066 per frame'' and you can see which version you would
rather debug at midnight.

\section{Timers}

The same logic drives anything that happens after a delay. Add \lstinline|dt| to
a counter and check it:

\begin{lstlisting}
float shootCooldown = 0.0f;

// update
if (shootCooldown > 0.0f) shootCooldown -= dt;

if (IsMouseButtonPressed(MOUSE_BUTTON_LEFT) && shootCooldown <= 0.0f)
{
    Shoot();
    shootCooldown = 0.25f;     // four shots per second, on any machine
}
\end{lstlisting}

Animation is the same pattern, and lesson 5 uses exactly this shape to step
through a sprite sheet.

\gotcha{Do not use \lstinline|GetTime()| for this. \lstinline|GetTime()| is
absolute seconds since startup: fine for a clock, awkward for ``0.25 seconds from
now'' because you have to remember when now was. Counting down a float with
\lstinline|dt| is simpler and pauses for free when you stop updating it.}

\section{The one case where delta time bites back}

If the game freezes for a second --- a big file loads, you dragged the window,
you hit a breakpoint --- the next \lstinline|dt| is huge. Everything moves a
second's worth in one frame, and fast objects teleport straight through walls,
because the collision check only looks at where they \emph{landed}, not the path
they took.

The standard defence is one line, and it is why you often see this in real games:

\begin{lstlisting}
float dt = GetFrameTime();
if (dt > 0.05f) dt = 0.05f;    // never simulate more than 50 ms in one step
\end{lstlisting}

The game slows down briefly instead of exploding. Good enough for everything you
will write for a long while.

\tryit{Lesson 3 puts a green square that uses \lstinline|dt| next to a red square
that does not, and drives both from the same keys.}{./course2d/build.sh 3}
